\subsection{High-dimensional experiments with QuadraSHAP}
\label{sec:high-dimensional-quadrashap}

To assess whether Shapley attribution remains practical when the number of
features greatly exceeds the number of patients, we consider two genomic
survival datasets: GSE24080 multiple-myeloma gene expression
\cite{geo_gse24080} and TCGA lower-grade glioma (LGG) DNA methylation
\cite{tcga_lgg_methylation}.
Genome-wide assays measure many molecular variables per patient, whereas
the cohorts with matched survival follow-up contain only hundreds of patients.
After eligibility filtering, GSE24080 contains 553 patients and 54,675
expression probe sets, and TCGA LGG contains 511 patients and 396,065
methylation sites. These are probe- or site-level features, retained without
gene aggregation. We use the supplied 339/214 training/validation split
for GSE24080 and a 383/128 split stratified by event status for TCGA LGG
(seed 42), giving an extreme small-$n$, large-$d$ setting.

We model overall survival using a Cox proportional-hazards model, which
incorporates right-censored observations: patients whose death has not been
observed by the end of follow-up. Its relative hazard factors across features,
making QuadraSHAP applicable directly:
\[
h(t\mid x)=h_0(t)f(x),
\qquad
f(x)=\exp(\beta^\top x)=\prod_{j=1}^{d}\exp(\beta_jx_j).
\]
To stabilize estimation when $d\gg n$, we shrink the Cox coefficients
towards zero using a ridge penalty, with the implementation's penalty
parameter fixed a priori at $\alpha=0.01d$. We use the Breslow approximation
for patients with equal recorded event times. Missing values are imputed using
training medians, and centering and scaling use training statistics.
All retained features have nonzero fitted coefficients.

We explain the relative hazard using 30 training-background patients,
sampled without replacement with seed 42 and held fixed across methods.
Specifically, absent features are replaced by the corresponding entries
of each background row and predictions are averaged:
\[
v_x(S)=\frac{1}{30}\sum_{b=1}^{30}
f\bigl(x_S,z^{(b)}_{\overline S}\bigr).
\]
The resulting game is a weighted sum of product games. We explain the
first five held-out patients in each saved split with the degree-exact rule
$m_q=\lceil d/2\rceil$ and with patient-specific node counts selected
to certify absolute quadrature tolerances
$\varepsilon\in\{10^{-1},10^{-2},10^{-3}\}$.
The timing experiments use JAX-Metal FP32 on an Apple M4 Pro, an 8\,GiB memory-planning
budget, and one background row per block. We record one synchronized
explanation call per patient and method after warming up each distinct
dataset/node-count configuration. Reported times include automatic node
selection, data transfers, and host accumulation, but exclude model fitting,
warm-up, and cached quadrature-rule construction.

The approximate GPU experiments were rerun for these three tolerances.
The degree-exact baseline is reused from the matching five-patient experiment;
the model, background rows, patient inputs, numerical source hashes,
software versions, and GPU memory plan were verified to match.

Table~\ref{tab:cox-high-dimensional} shows that degree-exact quadrature remains
manageable even at these dimensions: 27,338 and 198,033 nodes require
1.17\,s and 66.38\,s per patient on average, respectively.
For latency-sensitive use, the certified node budgets reduce the measured
latencies to 85--90\,ms for GSE24080
and 310--320\,ms for TCGA LGG.
These subsecond latencies support interactive, on-demand Shapley
explanations once the model and quadrature rules are ready. At
$\varepsilon=10^{-1}$, the speedups relative to the reused exact means are
$13.1\times$ and
$207.1\times$, respectively.
The exact-rule construction costs are substantial but can be amortized
across patients by caching the nodes and weights.

The analytical certificate and measured numerical discrepancy answer
different questions. For patient $p$, let $\phi_p$ be the exact Shapley
vector for the fixed background game and $\phi_{m_p,p}$ its quadrature
approximation in real arithmetic. The selected node count satisfies
\[
\|\phi_{m_p,p}-\phi_p\|_\infty\le C_p\le\varepsilon,
\qquad C_{\max}=\max_{p=1,\ldots,5}C_p.
\]
Thus $C_{\max}$ is an evaluated analytical bound, not an observed error.
The table's empirical discrepancy instead compares two floating-point outputs:
\[
\Delta_{\mathrm{GPU}}=
\max_p\|\widehat\phi^{\mathrm{GPU}}_{m_p,p}
-\widehat\phi^{\mathrm{GPU}}_{\mathrm{exact},p}\|_\infty.
\]
Every selected rule satisfies its analytical quadrature bound.
At $\varepsilon=10^{-1}$, the observed GPU
discrepancy is within the requested threshold for 5/5 GSE24080
and 5/5 TCGA LGG patients.
At $\varepsilon=10^{-2}$, the observed GPU
discrepancy is within the requested threshold for 5/5 GSE24080
and 5/5 TCGA LGG patients.
At $\varepsilon=10^{-3}$, the observed GPU
discrepancy is within the requested threshold for 4/5 GSE24080
and 2/5 TCGA LGG patients.

Agreement with the GPU degree-exact result alone does not establish exact
numerical accuracy, because that baseline also contains FP32 rounding error.
We therefore additionally compare both the GPU approximation and a separate
FP64 CPU evaluation at the same selected nodes with an independent,
tightly bounded 64-node FP64 reference. This reference has analytical
quadrature bounds below $10^{-10}$; a 96-node cross-check is also available
for the first patient of each cohort. It is a high-accuracy numerical
reference, rather than a full degree-exact FP64 computation.
The FP64 CPU evaluation is an accuracy diagnostic; its errors are never
paired with the GPU timings.

\begin{table}[t]
\centering
\small
\setlength{\tabcolsep}{4pt}
\begin{tabular}{llrrrr}
\toprule
 & & \multicolumn{2}{c}{GPU FP32} & \multicolumn{2}{c}{CPU FP64} \\
\cmidrule(lr){3-4}\cmidrule(lr){5-6}
Dataset & Tolerance & Worst error & Pass & Worst error & Pass \\
\midrule
GSE24080 & $\varepsilon=10^{-1}$ & $2.03\times10^{-3}$ & 5/5 & $1.21\times10^{-11}$ & 5/5 \\
GSE24080 & $\varepsilon=10^{-2}$ & $2.15\times10^{-3}$ & 5/5 & $1.28\times10^{-11}$ & 5/5 \\
GSE24080 & $\varepsilon=10^{-3}$ & $2.15\times10^{-3}$ & 1/5 & $1.28\times10^{-11}$ & 5/5 \\
TCGA LGG & $\varepsilon=10^{-1}$ & $3.05\times10^{-2}$ & 5/5 & $1.40\times10^{-10}$ & 5/5 \\
TCGA LGG & $\varepsilon=10^{-2}$ & $3.52\times10^{-2}$ & 3/5 & $1.31\times10^{-10}$ & 5/5 \\
TCGA LGG & $\varepsilon=10^{-3}$ & $3.91\times10^{-2}$ & 0/5 & $1.59\times10^{-10}$ & 5/5 \\
\bottomrule
\end{tabular}
\caption{Separate precision checks against the same independent FP64
reference. Worst error is the maximum absolute feature-wise error over all
five patients; Pass counts patients whose maximum error is at most their
requested tolerance. CPU FP64 entries are accuracy diagnostics and are
not the errors of the GPU runs timed in Table~\ref{tab:cox-high-dimensional}.}
\label{tab:cox-precision-checks}
\end{table}

The separate FP64 evaluation meets its requested threshold in 30/30
patient--tolerance cases, with a worst absolute difference of $1.59\times10^{-10}$,
whereas the timed FP32 GPU evaluations meet it in 19/30 cases
against that same reference.
For TCGA LGG methylation at $\varepsilon=10^{-2}$,
all five GPU approximations agree with their GPU degree-exact outputs
within the threshold, but only 3/5 agree with the independent
FP64 reference within that threshold (worst error $3.52\times10^{-2}$).
This illustrates why agreement between two FP32 computations alone
cannot establish the requested numerical accuracy.
These comparisons distinguish certified quadrature approximation from
end-to-end floating-point accuracy. Higher precision can reduce numerical
error, but the analytical certificate itself does not include rounding.
It also does not bound uncertainty from replacing a population background
distribution with the chosen 30-row empirical background.
